% =====================================================================================
%  Sec. 3 -- THE TWO BUDGETS.  The central theorem of the paper.  WRITTEN 2026-09-18.
%
%  This file did not exist until now.  thm:leakage, sec:leakage and eq:twobudget were
%  referenced 37, 17 and 7 times across the other nine sections and were defined
%  nowhere; app:counterex was referenced three times and defined nowhere.  All four
%  are defined here.
%
%  LABELS DEFINED HERE (and nowhere else):
%      sec:leakage  eq:lehmannA  eq:ritz3  eq:lambda  eq:twobudget  thm:leakage
%      eq:split  app:counterex
%  Sec. 3.1 (sec_3_1.tex, the retraction) is a SUBSECTION of this section and is
%  \input at the end of the [BODY] block.
%
%  PROVENANCE.  The statement, the three hypotheses and the numerical audit are those
%  adjudicated in P2_FRONTERA.md Sec. 4; the proof below is the one the evaluator
%  release/src/leakage_certificate.py realises, step for step, and its four rungs are
%  the four rungs whose contributions Sec. 5.5 measures (44 / 38 / 11 / 4.6 per cent).
%  The closed form of Lambda and the kernel 2 i eta / (Delta + 2 i eta) were re-derived
%  independently for this file by contour integration and agree with the form the code
%  evaluates.
%
%  VERIFIED FOR THIS FILE, on the 378 rows of the three deposited certificate scans
%  (release/data/cert_akw.json, cert_stress.json, cert_teqsci.json), by a script that
%  shares nothing with the one that wrote them:
%      - lower bound holds on 378 of 378;
%      - upper bound holds on 378 of 378; after removing the 2 rows that reproduce the
%        target to double-precision round-off (relative L1 <= 1e-10, where no ratio is
%        defined) the largest ratio error/bound is 0.9999999997;
%      - the exclusion lemma K_S >= 0 <=> w_S = 1 holds on 378 of 378, with 357 rows
%        at K_S = -1 and 21 at w_S = 1 exactly;
%      - the retracted bound of Sec. 3.1 fails on 219 of the 219 finite-ratio rows with
%        w_S >= 0.99, mildest factor 2.00, plus 19 rows where it is identically zero.
%
%  RESOLVED 2026-09-18: lowdin1962, feshbach1958 and saad2011 are in the bibliography and
%  were checked against Crossref (saad2011 is the SIAM *revised* edition, not a "2nd ed.").
%  This section additionally cites, since the novelty paragraph was rewritten:
%  paige1976, beattie2005, temple1928, weinstein1934, kato1949, frommersimoncini2008,
%  jialv2015, frommerschweitzer2016, oliveira2026 -- all present, all verified.
% =====================================================================================

% ======================================= [BODY] ======================================

\section{Two budgets: the missed weight and the boundary}
\label{sec:leakage}

The object a sample-based calculation delivers is a subspace, and the object a user wants is an error
bar on the response built inside it. This section proves the bound that connects them. It is two-sided,
it holds for \emph{any} orthogonal projector and therefore for any set of configurations a sampler
might return, and its upper branch separates into two terms with different economics: the first is fixed by the captured weight and vanishes once the probe's support is held; the second is not controlled by the weight, although it falls as the subspace grows. Everything after it is an evaluation of that
separation --- Sec.~\ref{sec:frontier} measures where the resulting certificate says something,
Sec.~\ref{sec:nogo} shows that no invariant of the state can replace it, and Sec.~\ref{sec:prices}
prices both terms in shots.

\subsection{The objects}

Fix the $(N{\pm}1)$ sector, let $H$ be the Hamiltonian restricted to it, $E_0$ the ground-state energy
of the $N$-electron sector, and let $\phi$ be the probe --- $\hat c^{\dagger}_p\ket{\Psi_0}$ or
$\hat c_p\ket{\Psi_0}$ for the charged channels, $\hat n_q\ket{\Psi_0}$ or $\hat S^{z}_q\ket{\Psi_0}$
for the neutral ones. Write $z=\omega+E_0+i\eta$ and $R(z)=(z-H)^{-1}$. The $\eta$-broadened spectral
function of Eq.~\eqref{eq:green} is
\begin{widetext}
\begin{equation}
A(\omega)\;=\;\frac{\eta}{\pi}\,\bigl\lVert R(z)\phi\bigr\rVert^{2},
\qquad\text{so that}\qquad
\int_{\mathbb R}A(\omega)\,d\omega\;=\;\lVert\phi\rVert^{2}
\label{eq:lehmannA}
\end{equation}
\end{widetext}
for every $\eta>0$, the sum rule holding because the Poisson kernel integrates to one on each Lehmann
pole. Let $P$ be an orthogonal projector on the sector, $Q=I-P$, and
\begin{widetext}
\begin{equation}
\phi_P=P\phi,\quad \phi_Q=Q\phi,\quad
w_P=\frac{\lVert\phi_P\rVert^{2}}{\lVert\phi\rVert^{2}},\quad
H_P=PHP,\quad R_P(z)=(z-H_P)^{-1}\big|_{\operatorname{ran}P},
\label{eq:ritz3}
\end{equation}
\end{widetext}
with Ritz pairs $H_P\ket{\tilde m}=\tilde E_m\ket{\tilde m}$ and $c_m=\braket{\tilde m|\phi}$. The
reconstruction of Eq.~\eqref{eq:ritz} is $A_P(\omega)=(\eta/\pi)\lVert R_P(z)\phi_P\rVert^{2}$, and it
obeys the same sum rule with $\lVert\phi_P\rVert^{2}$ on the right. When $P=P_{\mathcal S}$ is diagonal
in the determinant basis --- the only case a device produces --- we write $S$ for the configuration set
and $w_S$, $\Lambda_S$ for the quantities below.

The magnitude the theorem is indexed on is the Hamiltonian coupling that crosses the boundary of the
subspace, averaged over the resolution window rather than over the spectrum. With
$g_m=QH\ket{\tilde m}$,
\begin{widetext}
\begin{equation}
\Lambda_P(\eta)^{2}\;:=\;\frac{\eta}{\pi}\int_{\mathbb R}
\bigl\lVert QH\,R_P(z)\,\phi_P\bigr\rVert^{2}d\omega
\;=\;\sum_{m,m'}c_m\bar c_{m'}\,\braket{g_{m'},g_m}\,
\frac{2i\eta}{(\tilde E_{m'}-\tilde E_m)+2i\eta},
\label{eq:lambda}
\end{equation}
\end{widetext}
the second equality being exact and proved below; the sum is real, the $(m,m')$ and $(m',m)$ terms
being complex conjugates. Two features of Eq.~\eqref{eq:lambda} carry the practical content of this
paper. It is assembled from objects the reconstruction has already produced --- the Ritz vectors and
their coefficients --- plus one matrix--vector product $H\ket{\tilde m}$ per Ritz vector, so evaluating
it costs a small multiple of the diagonalization that produced $A_P$ and \emph{does not require} $A$.
And it vanishes exactly when $\operatorname{span}(S)$ contains $\mathcal K(H,\phi_P)$, the Krylov space of the retained part of the probe: $QHR_P(z)\phi_P$ vanishes for every $z$ exactly when $\mathcal K(H_P,\phi_P)$ is invariant under $H$, and it then coincides with $\mathcal K(H,\phi_P)$. For a subspace that holds the whole probe, $w_P=1$, the condition is that $\operatorname{span}(S)$ contain $\mathcal K(H,\phi)$ itself---far more demanding than it looks, as Sec.~\ref{sec:frontier:obstruction} measures.

\subsection{The bound}

\begin{theorem}[a posteriori leakage certificate]
\label{thm:leakage}
Let $H=H^{\dagger}$ act on the $(N{\pm}1)$ sector, let $\eta>0$, let $P$ be \emph{any} orthogonal
projector on that sector, and let $A$, $A_P$, $w_P$ and $\Lambda_P(\eta)$ be as in
Eqs.~\eqref{eq:lehmannA}--\eqref{eq:lambda}. Then
\begin{widetext}
\begin{equation}
\lVert\phi\rVert^{2}\bigl(1-w_P\bigr)
\;\le\;\lVert A-A_P\rVert_{1}\;\le\;
\min\Bigl\{\,\lVert\phi\rVert^{2}\bigl(1+w_P\bigr),\;
\bigl(\lVert\phi\rVert+\lVert\phi_P\rVert\bigr)
\bigl(\,\underbrace{\lVert\phi_Q\rVert}_{\text{missed weight}}
+\underbrace{\Lambda_P(\eta)/\eta}_{\text{boundary leakage}}\,\bigr)\Bigr\},
\label{eq:twobudget}
\end{equation}
% The restatement below used to sit in a SECOND widetext, with this clause as column
% text between the two.  That gave the output routine a legal break inside the theorem,
% and it took it: the hypothesis and Eq. (7) closed one page while the normalised form
% appeared on the next, behind Fig. 2 and its ten-line caption.  One widetext, no break
% point -- and two fewer horizontal rules on the page.
\noindent equivalently, dividing by $\lVert\phi\rVert^{2}$ and writing
$\widehat\Lambda_P=\Lambda_P/\lVert\phi\rVert$,
\begin{equation*}
1-w_P\;\le\;\frac{\lVert A-A_P\rVert_{1}}{\lVert\phi\rVert^{2}}\;\le\;
\min\Bigl\{\,1+w_P,\;\bigl(1+\sqrt{w_P}\bigr)
\bigl(\sqrt{1-w_P}+\widehat\Lambda_P(\eta)/\eta\bigr)\Bigr\}.
\end{equation*}
\end{widetext}
\end{theorem}


\noindent Three hypotheses carry it, each a property of the \emph{pipeline} that produced $A_P$ rather
than of the inequality: \emph{(H1)} $A$ and $A_P$ are referred to the same $E_0$; \emph{(H2)} the $L_1$
norm is taken over the whole line and normalized by $\lVert\phi\rVert^{2}$; \emph{(H3)} $A$ is the exact
spectral function of the full sector. Each is load-bearing and each is checked against this paper's own
pipeline in Sec.~\ref{app:hyp}; the proof, from the Feshbach--L\"owdin split of the resolvent
(Eq.~\eqref{eq:split}) and Cauchy--Schwarz in $\omega$, is in Sec.~\ref{sm:leakage}.

\subsection{What the two branches are, and what they are not}

\paragraph{What the weight does not control.} The first summand of the leakage branch, $\lVert\phi_Q\rVert=\lVert\phi\rVert\sqrt{1-w_P}$, is the Born weight of the probe the subspace failed to hold: it vanishes once the sampler holds the probe's Born support, and Table~\ref{tab:provenance} prices that purchase channel by channel. The second, $\Lambda_P(\eta)/\eta$, reads the off-diagonal block $PHQ$, and the captured weight does not control it: along the sampler's ranking it falls as $S$ grows (Sec.~\ref{sec:frontier}), but a subspace can hold the whole probe and still leak. Section~\ref{sec:retraction} turns this from a reading into a measurement: on
$\mathcal S^{\star}=\operatorname{supp}(\phi)$, where the first summand vanishes identically and by
construction, the relative $L_1$ error remains of order one, against a ceiling of $2$.

\paragraph{A certificate, not an error estimate.} The gap between the bound and the measured error is
large, and it is not the accident of a constant. At \emph{fixed} $P$ the looseness is stable --- within
$\pm17\%$ ($L=6$) and $\pm25\%$ ($L=8$) over a decade in $\eta$ --- but across subspaces it varies by a
factor $51$, and at the threshold where the leakage branch first improves on the trivial one it is
$3\times10^{2}$ to $2.4\times10^{3}$ (Sec.~\ref{sec:frontier:calib}). Eq.~\eqref{eq:twobudget} is therefore an exclusion, not an error
estimate; Sec.~\ref{sec:frontier} reports what it excludes, and at what fraction of the sector.

\paragraph{Only one of the two terms is a posteriori.} $\Lambda_P(\eta)$ is computable from $H$,
$\operatorname{ran}P$, $\phi$ and the Ritz vectors alone. The weight term is not: $w_P$ needs the exact
decomposition of the probe norm, and estimating it against an already-truncated $\Psi_0$ is biased in
the optimistic direction. The complete bound is therefore \emph{not} evaluable by a user who does not
already have the answer, and we never claim that it is; Sec.~\ref{sec:frontier:calib} gives the
conservative substitution that makes the decision itself a posteriori in the regime where it is used.

\paragraph{Prior art, and what is new.} Nothing in the proof is new as analysis. The split is L\"owdin
partitioning~\cite{lowdin1962}, the algebra of the Feshbach projection~\cite{feshbach1958}; Rayleigh--Ritz
and Krylov error analysis for resolvents and matrix
functions~\cite{saad2011,golub2010,epperly2022,temple1928,weinstein1934,kato1949,paige1976,beattie2005}
contains every rung of its steps, and the integration step is the reduced-basis bound on a Galerkin
projection with inf--sup constant $\eta$~\cite{rozza2008,hesthaven2016}. A posteriori estimates for Krylov
approximations to matrix functions are due to Frommer and Simoncini, Jia and Lv, and Frommer and
Schweitzer~\cite{frommersimoncini2008,jialv2015,frommerschweitzer2016}, with Lanczos and
stochastic-Lanczos-quadrature versions~\cite{chen2022lanczos,xuchen2025block,chen2021slq}; on quantum
Krylov, Oliveira, Ziems and Glaser assess reliability without knowledge of the true
spectrum~\cite{oliveira2026}; and a response-targeted selection criterion for classical selected
configuration interaction is Reinholdt, Kjellgren and Kongsted's~\cite{reinholdt2026lr}. The priority
for the a posteriori framing is theirs. What we add is the conjunction of three things: a closed form
for $\Lambda_P(\eta)$ from the Ritz pairs on a subspace of \emph{coordinates}, with no Krylov structure,
invariance or gap assumption; a bound two-sided on the same object; and a measured calibration of where
it crosses the trivial bound (Sec.~\ref{sec:frontier:calib}). Whether it is useful at the fractions
sampling reaches is answered, in the negative, in Sec.~\ref{sec:frontier}; the full comparison with that
literature is in Sec.~\ref{sm:leakage}.

% =====================================================================================
%  Sec. 3.1 -- Retraction of Theorem 1(iii).  Goes INSIDE Sec. 3 (Theorem B), in the body.
%  Supersedes: theorem_b2.tex lines 45-49 (the (iii) statement) and the (iii) clause of
%              the proof sketch, theorem_b2.tex lines 60-65;
%              method.tex lines 82-88 ("the genuine certificate ...");
%              discussion.tex lines 32-36 ("its error is bounded by the captured spectral weight").
%  Drop-in replacements for the latter two are appended, commented out, at the end of this file.
%
%  EXTERNAL LABELS USED HERE (must exist in the v3 tree -- see the report):
%    thm:leakage    Theorem B, the two-sided replacement            [Sec. 3, new; agreed]
%    sec:frontier   where the certificate's non-vacuity is mapped   [Sec. 5, new; agreed]
%    sec:moments    moment exactness, the surviving part (i)        [Sec. 4, defined in sec_4.tex]
%    sec:nogo       the orbital-invariance obstruction              [Sec. 6, defined in sec_6.tex]
%    app:counterex  counterexamples + the 292-subspace sweep        [Appendix B, NOT YET DEFINED]
%    fig:violation  new Fig. N3, violation vs fraction and vs eta   [Appendix B, NOT YET DEFINED]
%    sec:prices     the three prices                                [Sec. 7, defined in sec_7.tex]
%    tab:fracgrid   frac(L,eta) grid                                [Sec. 7, defined in sec_7.tex]
%    sec:method, sec:honest, fig:akwsampled, fig:scaling            [already in the tree]
%  This file defines only sec:retraction, eq:weightonly, eq:exclusion, eq:mu2.
%  Sec. 3's own label is sec:leakage -- agreed across sec_4.tex, sec_6.tex and sec_7.tex.
%  The obsolete alternative "sec:budgets" is used NOWHERE; do not reintroduce it.
%  NO \cite key exists for this preprint itself; arXiv:2608.16436 is given in prose.
%
%  %TODO(Sec. 3): three hypotheses of Theorem B are NOT stated in this file because they belong to
%  the theorem's own statement, which Sec. 3 owns.  They are load-bearing and a referee finds them by
%  reading step (1) of the proof.  They must appear in the statement, not in a reply:
%    (a) the published relative L1 is normalized by the in-window integral of |A_exact|, not by
%        ||phi||^2 (inflation ~1.027 at L=8);
%    (b) the "exact" reference of the resource scan is recomputed at the SAME n_l as the subspaces,
%        so the published relative L1 is ||A_K^full - A_K^S||, not ||A - A_K^S||;
%    (c) the theorem assumes A and A_P share E_0.  In the flagship point an origin error of
%        3.5e-4 t alone reproduces the whole certified error, and in a real shot-based pipeline E_0 is
%        itself estimated on a sampled subspace.  It has never been tested that way.
% =====================================================================================

\subsection{The captured weight is not the controlling magnitude}
\label{sec:retraction}

The certificate one would want for a sampled reconstruction indexes the error on the captured
weight alone,
\begin{equation}
\lVert A-A_S\rVert_{1}\;\le\;2\lVert\phi\rVert^{2}(1-w_S)\;+\;C\,\rho^{-K_S},
\label{eq:weightonly}
\end{equation}
with $w_S=\lVert P_S\phi\rVert^{2}/\lVert\phi\rVert^{2}$ and $K_S$ the largest Krylov order
contained in $\operatorname{span}(S)$ (we set $K_S=-1$ when $\phi\notin\operatorname{span}(S)$).
It is false, for every constant $C$ and every $\rho$, and it fails hardest in exactly the regime
a sampling experiment aims for. 

\paragraph{A two-level counterexample forces the trivial constant.}
Take $H=\big(\begin{smallmatrix}0&g\\ g&0\end{smallmatrix}\big)$, $\phi=e_0$, $S=\{0\}$, $E_0=0$. Then
$P_S\phi=\phi$, so $w_S=1$ exactly and the weight term vanishes identically; and
$\mathcal K_0=\operatorname{span}\{\phi\}\subseteq\operatorname{span}(S)$ while
$\mathcal K_1\not\subseteq\operatorname{span}(S)$, so $K_S=0$ exactly and the second term is the bare
constant $C$. The exact $A$ is a pair of Lorentzians of weight $\tfrac12$ at $\pm g$; the
reconstruction, built on $H_S=P_SHP_S=0$, is one Lorentzian of weight $1$ at the origin, so
$\lVert A-A_S\rVert_{1}\to2\lVert\phi\rVert^{2}$ as $g/\eta\to\infty$. Since $g$ is free,
\eqref{eq:weightonly} requires $C\ge2\lVert\phi\rVert^{2}$: false if $C$ is a constant, vacuous if it may depend on the instance. The same holds at every Krylov order: on a chain of $K+2$ sites with $\phi=e_0$ and $S$ its first $K+1$ sites, $w_S=1$ and $K_S=K$, and the exact and reconstructed spectra have strictly interlacing poles, so the same limit requires $C\rho^{-K}\ge2\lVert\phi\rVert^{2}$ for every $K$, which no fixed $C$ satisfies when $\rho>1$. With $w_S<1$ strictly the weight term alone fails without bound: in a three-level family its violation diverges as $\varepsilon^{-2}$ in the amplitude $\varepsilon$ of the discarded component (Sec.~\ref{app:counterex}). By contrast
$\lVert A-A_S\rVert_{1}\le\lVert\phi\rVert^{2}(1+w_S)$ holds for every $S$ under no hypothesis at all,
$A$ and $A_S$ being nonnegative with known integrals: that trivial bound is the only surviving member of
the family \eqref{eq:weightonly}, and Theorem~\ref{thm:leakage} keeps it explicitly, as the first branch
of its minimum.
\paragraph{The violation is systematic, and worst where the reconstruction is best.} In a pooled sweep
of $606$ subspaces from five Hamiltonian families the weight-only bound fails on $468$, and on all $363$
with $w_S\ge0.99$, the mildest by a factor $2.10$ (Fig.~\ref{fig:thm1iii-violation}); in a
first sweep it failed on $126$ of $126$ such subspaces. In the
configuration of Fig.~\ref{fig:akwsampled} its median violation factor rises from $18.5$ to $29.9$
between sector fractions $0.50$ and $0.95$ while the measured error falls by three orders of magnitude.
The sweep and its declared limits are in Sec.~\ref{sm:leakage}; the conclusion does not depend on it,
the counterexample being analytic.

\paragraph{Why it fails.} The reconstruction is Rayleigh--Ritz on $H_S=P_SHP_S$, not a restriction of
$A$ to a subset of its eigenpairs: weight the subspace keeps is \emph{displaced} in $\omega$, not
discarded, and the two-level example is the extreme case. What the integrals do give is the lower bound
$\lVert\phi\rVert^{2}(1-w_S)\le\lVert A-A_S\rVert_{1}$, with $0$ violations across the $606$ subspaces,
the left-hand half of Theorem~\ref{thm:leakage}.

\paragraph{The two terms are mutually exclusive, and one is off at every operating point here.}
$S$ is a set of determinants, so $P_S$ is diagonal in the working basis and
$P_S\phi=\phi\iff\operatorname{supp}(\phi)\subseteq S\iff w_S=1$, while $K_S\ge0$ means precisely
$\phi=H^{0}\phi\in\operatorname{span}(S)$. Hence the exclusion lemma
\begin{equation}
K_S\ \ge\ 0\qquad\Longleftrightarrow\qquad w_S\ =\ 1 ,
\label{eq:exclusion}
\end{equation}
verified with $0$ discrepancies across the $378$ deposited subspaces, of which $357$ have $K_S=-1$ (Sec.~\ref{app:counterex}). The two terms of
\eqref{eq:weightonly} are therefore never both informative: either $w_S=1$ and the whole bound is the
second term --- the two-level counterexample --- or $K_S=-1$ and $C\rho^{-K_S}=C\rho$, an undefined
constant with no decay in it. Counting settles which branch applies here. The seed is
$\phi=\hat c^{\dagger}_{0\uparrow}\ket{\Psi_0}$ at half filling, so $\operatorname{supp}(\phi)$ lies in the set of $(N{+}1)$ determinants with orbital $0\!\uparrow$ occupied, a fraction $n_\uparrow/L=\tfrac12+\tfrac1L$ of the sector, and fills it except where symmetry zeroes an amplitude of $\Psi_0$. By \eqref{eq:exclusion}, $K_S\ge0$ demands $|S|\ge\lvert\operatorname{supp}(\phi)\rvert$, a floor of $0.571$ at $L=8$ ($2237$ of $3920$), $0.600$ at $L=10$, and above $0.54$ at $L=12$ and $14$ at any amplitude threshold from $10^{-12}$ to $10^{-10}$, against the sector fractions $0.5561$, $0.3625$, $0.1700$, $0.0800$
at which the resource scan of Sec.~\ref{sec:prices} operates (Fig.~\ref{fig:scaling}; the coarser
replication grid of Table~\ref{tab:fracgrid} reads $0.180$ at $L=12$ for the same operating point).
From $L=8$ on, every operating point of that scan lies below its own floor, so $K_S=-1$ by counting
alone, with no
numerics, no ranking and no protocol, and any point below the same line inherits the verdict
(Sec.~\ref{app:counterex}). 
Theorem~\ref{thm:leakage} carries no such
bookkeeping and is held to the opposite standard. Its two terms are defined and simultaneously active
for every $S$; the weight term is the half this work cannot compute without $\Psi_0$ ---
$\lVert\phi_Q\rVert=\lVert\phi\rVert\sqrt{1-w_S}$, and estimating $w_S$ against an already-truncated
$\Psi_0$ is biased optimistically --- while $\Lambda_S(\eta)$ is an \emph{a posteriori} quantity
assembled from the Ritz vectors the reconstruction already forms plus one matrix--vector product, and is
evaluated, subspace by subspace, in Sec.~\ref{sec:frontier}.

\paragraph{The magnitude, not the constant.}
Two measurements at the same, exactly unit, captured weight show that \eqref{eq:weightonly} was indexed
not on a bad constant but on the wrong quantity. First, let $S^\star=\operatorname{supp}(\phi)$, the
smallest determinant subspace of unit captured weight --- fixed by the seed, with no ranking, no
tie-breaking and no zero-amplitude padding --- of size $(\tfrac12+\tfrac1L)D$ at $L=6$ and $0.571\,D$ at $L=8$, where $426$ of the $4900$ ground-state amplitudes vanish by symmetry. Its weight term
is identically zero, and its relative $L_1$ error is $0.43$ at $L=6$ and $0.35$ at $L=8$ on the periodic rings at $\eta=0.18t$, rising to $1.15$ and $0.73$ at $\eta=0.05t$, against a hard ceiling of $1+w_S=2$.
Second, $S^\star$ is the $K=0$ rung of the Krylov ladder
$S_K=\bigcup_{a\le K}\operatorname{supp}(H^{a}\phi)$, every rung of which contains $\phi$ and so carries
$w_S=1$ identically, while its moment reach grows: $S_K$ reproduces the exact spectral moments through
order $2K{+}1$ --- to the $10^{-15}$ level, failing as it must at $2K{+}2$ --- whereas coordinate
subspaces of the same size drawn at random miss already $\mu_0$, in all of $300$ draws
(Sec.~\ref{sec:moments}). Along a family of constant, exactly unit captured weight the error therefore
runs from order unity to moment-exact through order $2K{+}1$. The captured weight, identical on every
rung, does not order them.

What does is elementary. Writing $\mu_j,\mu^{S}_j$ for the moments of the unbroadened measures of
Sec.~\ref{sec:moments}, every determinant subspace with $w_S=1$ reproduces $\mu_0$ and $\mu_1$ exactly,
and the first moment that can fail is $\mu_2$, with deficit
\begin{equation}
\mu_2-\mu^{S}_2\;=\;\lVert Q_S H P_S\phi\rVert^{2},\qquad Q_S=I-P_S ,
\label{eq:mu2}
\end{equation}
the squared coupling of the Hamiltonian across the boundary of $S$ ($E_0$ drops out because
$Q_S\phi=0$), zero exactly when $S$ also holds $\operatorname{supp}(H\phi)$ --- the $K_S\ge1$ condition.
That boundary coupling is the magnitude the replacement is indexed on: it controls the leakage term
through $\Lambda_S(\eta)\le\lVert P_SHQ_S\rVert\,\lVert\phi_S\rVert$, and Theorem~\ref{thm:leakage}
averages it over the resolution window rather than over the spectrum. The captured weight says what the
subspace \emph{contains}; the error is set by what the Hamiltonian does to the seed at the subspace's
\emph{edge}. Both halves of this paper follow from that displacement: Sec.~\ref{sec:frontier} measures
where the boundary term is small enough for the certificate to say something, and Sec.~\ref{sec:nogo}
shows why no orbital-rotation invariant can predict it in advance.

% =====================================================================================
%  DROP-IN REPLACEMENTS FOR THE TWO SURVIVING INVOCATIONS (not part of Sec. 3.1;
%  move into the target files -- do not compile from here).
%
%  (1) method.tex, replacing the sentence at lines 82-88 ("These zeroth-moment checks are not
%      themselves an error bound; the genuine certificate is ... (Sec.~\ref{sec:honest})."):
%
% These zeroth-moment checks are not themselves an error bound. The certificate is the two-sided
% bound of Theorem~\ref{thm:leakage}, which splits the $L_1$ error into the Born weight the subspace
% failed to capture and the Hamiltonian coupling leaking across its boundary at resolution $\eta$;
% only the first of the two is bought by drawing more bitstrings. A weight-only bound appeared in
% versions~1 and~2 of this preprint and is retracted in Sec.~\ref{sec:retraction}. Where the
% certificate is non-vacuous is measured, not asserted, in Sec.~\ref{sec:frontier} --- an essential
% property given that the classical hardness of a specific instance is empirical rather than a
% theorem (Sec.~\ref{sec:honest}).
%
%  (2) discussion.tex, replacing lines 32-36 ("On the positive side, Theorem~\ref{thm:b2}
%      (Sec.~\ref{app:b2}) guarantees ... bounded by the captured spectral weight."):
%
% On the positive side, moment exactness (Sec.~\ref{sec:moments}) guarantees that a subspace
% containing the order-$K$ Krylov space reproduces the exact spectral moments through order $2K{+}1$
% and is captured with a polynomial shot budget, and Theorem~\ref{thm:leakage} certifies the
% reconstruction two-sidedly: from below by the captured weight, from above by that weight together
% with the boundary leakage. The error is \emph{not} bounded by the captured spectral weight alone ---
% a subspace holding the entire Born support of the seed, of unit weight by construction, still
% carries a relative $L_1$ error of $0.28$ to $0.43$ at $\eta=0.18t$, the lower end on the open chain
% of the method validation and the upper on the periodic ring (Sec.~\ref{sec:retraction}).
%
%  INTEGRATION NOTE: after these two replacements are made, \ref{thm:b2} must appear NOWHERE in the
%  live tree.  Its parts (i) and (ii) become Theorem~\ref{thm:moments} and Prop.~\ref{prop:shots}
%  (Sec. 4), its part (iii) is retracted here, and its corollary becomes Conjecture~\ref{conj:geom}.
%  The only occurrence of the string thm:b2 in this file is inside the quotation above, which is the
%  OLD sentence being replaced.
% =====================================================================================



